% =========================================================================
% NPL core paper — draft v0.5 (2026-07-03): round-2 fixes (cref slip in Thm 6.4(iii),
% collapse-theorem attribution via secondary, single label, opposite signs, abstract
% trim) + post-2022 sweep engaged (bilkova2021constraint, villadsen2017formalizing)
% + Prop 9.1(iii) inexpressibility of negative signs (Lean-verified).
% v0.4 (superseded): full revision per panel review round 1
% papers/npl-core/reviews/2026-07-03-referee-reports.md. Changes:
%  [1] stale compactness remark deleted; bifilter comparison RESOLVED (Prop 9.1);
%  [2] Contribution 1 rewritten (unsigned subsumed; negative signs provably
%      outside the bifilter framework; tau-invariance stated honestly);
%  [3] full 16-rule tableau table + explicit ND rule list printed;
%  [4] \models_FOUR defined; indicator notation defined (\ind, \mathbf{1});
%  [5] full compactness proof printed (non-Lean at the time; Lean-verified
%      2026-07-04, marker upgraded to [L] — see Thm 6.4);
%  [6] proof sketches added to all [L]-only results; [L] coverage made accurate;
%  [7] artifact statement (toolchain, build, coverage policy);
%  [8] draft residue removed; bibliographystyle plain; sect. 7 retitled;
%  [9] Fitting engagement, neologism footnote, tau-invariance (Prop 3.4).
% Build: pdflatex main && bibtex main && pdflatex main && pdflatex main
% =========================================================================
\documentclass[11pt]{article}

\usepackage[T1]{fontenc}
\usepackage[utf8]{inputenc}
\usepackage{lmodern}
\usepackage{amsmath,amssymb,amsthm}
\usepackage{booktabs,array}
\usepackage[colorlinks=true,linkcolor=blue,citecolor=blue,urlcolor=blue]{hyperref}
\usepackage[capitalize]{cleveref}
\usepackage{microtype}
% Allow long \texttt Lean identifiers to break at their (escaped) underscores.
\newcommand{\lb}{\discretionary{}{}{}}
\emergencystretch=3em

\newtheorem{theorem}{Theorem}[section]
\newtheorem{lemma}[theorem]{Lemma}
\newtheorem{proposition}[theorem]{Proposition}
\newtheorem{corollary}[theorem]{Corollary}
\theoremstyle{definition}
\newtheorem{definition}[theorem]{Definition}
\newtheorem{remark}[theorem]{Remark}

% ----- notation -----
\newcommand{\Frm}{\mathrm{Form}}
\newcommand{\FOUR}{\mathbf{4}}
\newcommand{\vT}{\mathbf{T}}
\newcommand{\vF}{\mathbf{F}}
\newcommand{\vB}{\mathbf{B}}
\newcommand{\vN}{\mathbf{N}}
\newcommand{\hmz}{\oplus}                       % harmonization (consensus)
\newcommand{\lek}{\le_{k}}
\newcommand{\tle}{\le_{t}}
\newcommand{\proj}[1]{\pi_{#1}}
\newcommand{\ind}[1]{\mathbf{1}[#1]}            % indicator
\newcommand{\Tp}{T^{+}}\newcommand{\Tm}{T^{-}}
\newcommand{\Fp}{F^{+}}\newcommand{\Fm}{F^{-}}
\newcommand{\dershort}{\vdash_{\!A}}
\newcommand{\ndder}{\vdash_{\mathrm{ND}}}
\newcommand{\lean}[1]{\texttt{#1}}

\title{Nullivance Propositional Logic:\\
Threshold-Signed Consequence on the Unit Square\\
with a Machine-Checked Completeness Theorem\thanks{%
Draft v0.5 (July 2026). Results marked [L] are verified in the Lean~4 artifact
described in Appendix~\ref{app:lean}; results marked [paper] carry full proofs
here and are not yet formalized. Panel-review history and revision log are kept
with the source.}}
\author{Trinh Tung Lam\thanks{Contact: \texttt{lamyofficialvn@gmail.com}.}}
\date{July 2026 --- draft v0.5}

\begin{document}
\maketitle

\begin{abstract}
We present the propositional core of \emph{nullivance logic} (NPL). Formulas
take values in the unit square $[0,1]^2$ --- independent degrees of support
for truth and for falsity; negation swaps the channels, conjunction and
disjunction act as $(\min,\max)$ and $(\max,\min)$, and a
\emph{harmonization} connective $\hmz$ acts as $(\min,\min)$: bilattice
consensus as a first-class connective. A \emph{manifestation threshold}
$\tau\in(0,1]$ induces four sign predicates $\Tp,\Tm,\Fp,\Fm$; consequence
is defined uniformly in the threshold and shown threshold-invariant (a
well-definedness property, not an added strength). Main results, each machine-checked in Lean~4 unless noted:
(i)~thresholding commutes exactly with evaluation, projecting the continuous
semantics onto the four-valued matrix $\FOUR$; (ii)~a four-signed analytic
tableau calculus is sound and complete for signed consequence, over $\FOUR$
and, via the exact projection, over the continuous semantics; consequence is
decidable, compact and strongly complete for infinite
premise sets; (iii)~NPL is paraconsistent --- explosion, disjunctive
syllogism, and detachment for the material conditional all fail --- while
harmonized contradictions $\varphi\hmz\lnot\varphi$ collapse to an unmanifest
state; (iv)~$\hmz$ is not definable from the $\{\lnot,\wedge,\vee\}$ fragment,
and the natural-deduction rule set of the source system is provably
incomplete for it --- its rules constrain only the truth channel of $\hmz$;
(v)~a generative tier beneath the semantics derives truth-objects from
(intensity, structure) channel pairs through a replaceable stability
function, with the logic provably independent of that choice. We position
NPL against Belnap--Dunn logic, bilattice logics, and paraconsistent
G\"odel logic: the semantic ingredients are known, the unsigned fragment is
an instance of Arieli--Avron's logical-bilattice consequence, and the
contribution is the signed architecture (which provably falls outside the
bifilter framework, and is provably not internalizable in the object
language), its exact projection, and the verified metatheory.
\end{abstract}

\section{Introduction}\label{sec:intro}

Two ideas drive nullivance logic.\footnote{The name derives from the
system's founding concern with \emph{structured absence}: states that are
null in intensity yet carry direction (Latin \emph{nullus} + valence). The
source material calls the pre-logical version of this notion
\emph{quasivance}; both terms are defined where they are used
(\cref{sec:generative}) and are otherwise dispensable labels for the formal
objects.} First, the epistemic state of a proposition has two independent
coordinates: how strongly it is supported as true, and how strongly as
false. This is the two-channel reading familiar from Belnap--Dunn
four-valued logic \cite{belnap1977useful,belnap1977computer,
dunn1976intuitive} and its continuous relatives; we inherit that literature
rather than compete with it. Second, \emph{manifestation is thresholded}: a
proposition manifests as true (false) only when the corresponding support
clears a threshold $\tau$. Below both thresholds lies an unmanifest state;
above both lies manifest contradiction. A designated connective $\hmz$
(\emph{harmonization}) takes the consensus of both channels, and
harmonizing a proposition with its own negation drives it into the
unmanifest state instead of exploding.

\paragraph{Honest positioning.} None of our semantic ingredients is new in
isolation, and we say so precisely (\cref{sec:related}): the
$\{\lnot,\wedge,\vee\}$ fragment of the continuous semantics coincides with
the propositional base of paraconsistent G\"odel logic
\cite{bilkova2022paraconsistent} and with product-bilattice semantics
\cite{ginsberg1988multivalued,fitting1991bilattices}; $\hmz$ is the
bilattice consensus $\otimes$, which at the four-valued level carries
complete calculi \cite{arieli1996reasoning}; the projection core is the
cutworthiness of $\alpha$-cuts \cite{klir1995fuzzy}; signed tableaux for
finitely-valued logics are classical \cite{haehnle1994automated,
priest2008introduction}. Moreover --- a comparison we performed rather than
conjectured --- the \emph{unsigned} fragment of our consequence relation is
an instance of Arieli--Avron's logical-bilattice consequence
(\cref{prop:bifilter}). The contributions are:

\begin{enumerate}
\item \textbf{Signed consequence beyond designated values}
  (\cref{def:consequence}, \cref{prop:bifilter}): four sign predicates read
  the two channels against the threshold. The negative signs $\Tm,\Fm$ are
  complements, and we prove their satisfaction sets are not bifilters; the
  four-signed relation is therefore not a logical-bilattice matrix
  consequence in the sense of \cite{arieli1996reasoning}, while its unsigned
  fragment provably is. The relation is threshold-invariant
  (\cref{prop:tau}), which we state as a well-definedness result.
\item \textbf{Exact projection} (\cref{thm:projection}): thresholding
  commutes with evaluation, collapsing continuous consequence onto the
  finite matrix $\FOUR$ without loss.
\item \textbf{A machine-checked completeness theorem}
  (\cref{thm:complete-cont}), plus decidability (\cref{thm:decidable}) and
  compactness and strong completeness
  (\cref{thm:compactness}). To our knowledge no verified completeness
  theorem exists for a consensus-bearing logic over $[0,1]^2$ with signed
  consequence (search scope in \cref{sec:related}).
\item \textbf{A calibrated failure map} (\cref{sec:profile}): explosion,
  disjunctive syllogism, and detachment fail with machine-checked finite
  countermodels; the natural-deduction rules of the source system are
  derived rules of the calculus, yet the ND system they generate is
  \emph{provably incomplete} ($\hmz$-self-duality is valid but underivable
  --- \cref{prop:nd}); $\hmz$ is interderivable with $\wedge$ but not
  congruent to it, and not definable from the FDE fragment
  (\cref{prop:primitive}).
\item \textbf{A verified generative tier} (\cref{sec:generative}):
  truth-objects arise from (intensity, structure) channel pairs through an
  axiomatized, replaceable stability function; independence of the logic
  from that choice is witnessed by the artifact's import graph, and the
  limitative facts (unmanifest projection, non-injectivity of
  initialization) are machine-checked.
\end{enumerate}

Results carry [L:\,\lean{name}] markers naming the verifying Lean
declaration; results not covered by the artifact are marked [paper] and
carry full proofs here. The artifact policy is in
Appendix~\ref{app:lean}.

\section{Syntax}\label{sec:syntax}

\begin{definition}[Language]\label{def:lang}
Fix a countably infinite set of atoms $p_0,p_1,\dots$ (countability is used
in \cref{thm:compactness}). Formulas:
\[
\varphi \;::=\; p \mid \lnot\varphi \mid \varphi\wedge\varphi \mid
\varphi\vee\varphi \mid \varphi\hmz\varphi .
\]
$\varphi\Rightarrow\psi$ abbreviates $\lnot\varphi\vee\psi$; it is not a
primitive and does not detach (\cref{prop:mp-fails}).
[L:\,\lean{Nullivance.Syntax.Formula}]
\end{definition}

\section{Two-channel semantics and threshold signs}\label{sec:semantics}

\begin{definition}[Truth-objects; models; valuation]\label{def:model}
A \emph{truth-object} is a pair $(t,f)\in[0,1]^2$: independent degrees of
support for truth and falsity. An \emph{NPL model} is $M=(v,\tau)$ with $v$
assigning a truth-object to every atom and $\tau\in(0,1]$ the
\emph{manifestation threshold}. $v$ extends to
$V_M(\varphi)=(t_M(\varphi),f_M(\varphi))$ by
\[
\setlength{\arraycolsep}{2pt}
\begin{array}{ll}
V_M(\lnot\psi) = (f_M(\psi),\,t_M(\psi)), \quad &
V_M(\psi\wedge\chi) = (\min(t,t'),\,\max(f,f')),\\[2pt]
V_M(\psi\vee\chi) = (\max(t,t'),\,\min(f,f')), \quad &
V_M(\psi\hmz\chi) = (\min(t,t'),\,\min(f,f')).
\end{array}
\]
writing $(t,f)=V_M(\psi)$, $(t',f')=V_M(\chi)$. The
$\{\lnot,\wedge,\vee\}$ clauses are the twist-product/bilattice clauses of
\cite{bilkova2022paraconsistent,fitting1991bilattices}; $\hmz$ is consensus.
[L:\,\lean{Nullivance.Continuous.TruthObj}, \lean{evalC}]
\end{definition}

\begin{definition}[Signs; states; consequence]\label{def:consequence}
$M\models\Tp\varphi \iff t_M(\varphi)\ge\tau$;
$M\models\Tm\varphi \iff t_M(\varphi)<\tau$;
$M\models\Fp\varphi \iff f_M(\varphi)\ge\tau$;
$M\models\Fm\varphi \iff f_M(\varphi)<\tau$.
The two thresholds classify $\varphi$ into four \emph{states}: manifest true
$\vT$, manifest false $\vF$, manifest contradiction $\vB$ ($\Tp$ and $\Fp$),
unmanifest $\vN$ ($\Tm$ and $\Fm$). Unsigned satisfaction is $\Tp$;
designated states are $\{\vT,\vB\}$. For a set $\Sigma$ of signed formulas:
$\Sigma\models S\varphi$ iff every model (every valuation, every
$\tau\in(0,1]$) satisfying $\Sigma$ satisfies $S\varphi$.
[L:\,\lean{Nullivance.Continuous.SatC}, \lean{Metatheory.ConsequenceC}]
\end{definition}

\begin{proposition}[Threshold invariance]\label{prop:tau}
For every fixed $\tau_0\in(0,1]$, the relation ``every valuation at
threshold $\tau_0$ satisfying $\Sigma$ satisfies $S\varphi$'' coincides with
$\models$. Quantifying over thresholds in \cref{def:consequence} is thus a
well-definedness statement, not an added strength.
[L:\,\lean{ConsequenceCAt}, \lean{consequenceCAt\_iff\_consequenceC}]
\end{proposition}
\begin{proof}[Proof sketch]
Both lifting directions of \cref{sec:metatheory} are uniform in the
threshold: a $\tau_0$-model projects to a $\FOUR$-valuation satisfying the
same signed formulas, and a $\FOUR$-valuation embeds at the corners
$\{0,1\}^2$, where $\ind{1\ge\tau_0}=1$ and $\ind{0\ge\tau_0}=0$ for
\emph{every} $\tau_0\in(0,1]$. Hence each fixed-$\tau_0$ relation equals
$\FOUR$-consequence (\cref{def:four}), hence equals $\models$.
\end{proof}

\begin{lemma}[Structural facts]\label{lem:structure}
(i) $V_M(\varphi)\in[0,1]^2$ (induction; $[0,1]$ is closed under
$\min,\max$). (ii) \emph{Latent collapse}:
$V_M(\varphi\hmz\lnot\varphi)=(m,m)$ with
$m=\min(t_M(\varphi),f_M(\varphi))$ --- both channels of the harmonized
contradiction agree, so for $\tau>m$ it is in state $\vN$. (iii) $\hmz$ is
commutative, associative, idempotent, self-dual under $\lnot$, with unit
$(1,1)$ on the square. (iv) With
$(t_1,f_1)\tle(t_2,f_2)\iff t_1\le t_2 \wedge f_2\le f_1$ and
$(t_1,f_1)\lek(t_2,f_2)\iff t_1\le t_2\wedge f_1\le f_2$: $\wedge/\vee$ are
the $\tle$ meet/join, $\hmz$ is the $\lek$ meet, and $(0,0),(1,1)$ are the
$\lek$ extrema --- the square is a reduct of the product bilattice
$[0,1]\odot[0,1]$ \cite{ginsberg1988multivalued} with $\hmz$ its consensus;
the $\lek$-join (gullibility) is deliberately not a connective.
[L:\,\lean{eval\_mem}, \lean{latent\_collapse}, \lean{oplus2\_*},
\lean{le\_t}/\lean{le\_k} family]
\end{lemma}

\section{The four-valued matrix and exact projection}\label{sec:projection}

\begin{definition}[$\FOUR$; $\FOUR$-consequence; projection]\label{def:four}
$\FOUR=\{0,1\}^2\subseteq[0,1]^2$, corners $\vT=(1,0)$, $\vF=(0,1)$,
$\vB=(1,1)$, $\vN=(0,0)$, closed under the clauses of \cref{def:model};
the $\{\lnot,\wedge,\vee\}$ tables are the Belnap--Dunn tables with
designated $\{\vT,\vB\}$ \cite{belnap1977useful,dunn1976intuitive}, and
$\hmz$ restricts to consensus $\otimes$ \cite{arieli1996reasoning}. A
$\FOUR$-valuation $w$ maps atoms to $\FOUR$ and extends by the same
clauses; it satisfies $\Tp\varphi$ ($\Fp\varphi$) iff the first (second)
bit of $w(\varphi)$ is $1$, and the negative signs are the complements.
\emph{$\FOUR$-consequence} $\Sigma\models_{\FOUR}S\varphi$: every
$\FOUR$-valuation satisfying $\Sigma$ satisfies $S\varphi$. The
\emph{threshold projection} is
$\proj{\tau}(x,y)=(\ind{x\ge\tau},\ind{y\ge\tau})$, where $\ind{\cdot}$ is
the $\{0,1\}$-valued indicator.
[L:\,\lean{Nullivance.Semantics.V4}, \lean{Metatheory.Consequence4},
\lean{Nullivance.Continuous.proj}]
\end{definition}

\begin{theorem}[Exact projection]\label{thm:projection}
For every model $M=(v,\tau)$ and formula $\varphi$:
$\proj{\tau}(V_M(\varphi)) = V^{\FOUR}_{\proj{\tau}\circ v}(\varphi)$.
Consequently $M\models S\varphi$ iff the projected $\FOUR$-valuation
satisfies $S\varphi$. [L:\,\lean{exact\_projection}, \lean{sat\_projection}]
\end{theorem}
\begin{proof}[Proof sketch]
Structural induction. The binary cases are the facts
$\ind{\min(x,y)\ge\tau}=\min(\ind{x\ge\tau},\ind{y\ge\tau})$ and dually for
$\max$ --- cutworthiness of the standard fuzzy operations
\cite{klir1995fuzzy}; negation commutes with the two cuts because it swaps
coordinates. The content is the use: consequence reduces \emph{exactly} to
$\FOUR$, with no approximation.
\end{proof}

\section{The four-signed analytic calculus}\label{sec:calculus}

\begin{definition}[Calculus]\label{def:calculus}
A \emph{branch} is a finite set of signed formulas; \emph{closed} iff it
contains $S\varphi$ and $\bar S\varphi$ for some $\varphi$ (opposite signs:
$\overline{\Tp}=\Tm$, $\overline{\Fp}=\Fm$, and conversely
$\overline{\Tm}=\Tp$, $\overline{\Fm}=\Fp$). The sixteen decomposition
rules, one per sign--connective pair (comma = add both to the branch;
$\mid$ = split the branch):

\begin{center}
\begin{tabular}{l@{\qquad}llll}
\toprule
 & $\Tp$ & $\Tm$ & $\Fp$ & $\Fm$ \\
\midrule
$\lnot\varphi$ &
  $\Fp\varphi$ & $\Fm\varphi$ & $\Tp\varphi$ & $\Tm\varphi$ \\
$\varphi\wedge\psi$ &
  $\Tp\varphi,\Tp\psi$ & $\Tm\varphi\mid\Tm\psi$ &
  $\Fp\varphi\mid\Fp\psi$ & $\Fm\varphi,\Fm\psi$ \\
$\varphi\vee\psi$ &
  $\Tp\varphi\mid\Tp\psi$ & $\Tm\varphi,\Tm\psi$ &
  $\Fp\varphi,\Fp\psi$ & $\Fm\varphi\mid\Fm\psi$ \\
$\varphi\hmz\psi$ &
  $\Tp\varphi,\Tp\psi$ & $\Tm\varphi\mid\Tm\psi$ &
  $\Fp\varphi,\Fp\psi$ & $\Fm\varphi\mid\Fm\psi$ \\
\bottomrule
\end{tabular}
\end{center}

The $\hmz$ row: the $\Tp$ column agrees with $\wedge$ (both truth channels
are $\min$), while $\Fp(\varphi\hmz\psi)$ is \emph{conjunctive}, unlike the
branching $\Fp(\varphi\wedge\psi)$ --- this is where $\hmz$ and $\wedge$
part ways inside the proof system. $\Sigma\dershort S\varphi$ iff some
tableau for $\Sigma\cup\{\bar S\varphi\}$ closes. The four signs are
sets-as-signs \cite{haehnle1994automated}: $\Tp=\{\vT,\vB\}$,
$\Tm=\{\vF,\vN\}$, $\Fp=\{\vF,\vB\}$, $\Fm=\{\vT,\vN\}$.
[L:\,\lean{Nullivance.ProofTheory.Closes}: an inductive closability
predicate with two closure and sixteen rule constructors; branching rules
take two subderivations]
\end{definition}

\section{Soundness, completeness, decidability, compactness}\label{sec:metatheory}

\begin{theorem}[Soundness and completeness over $\FOUR$]\label{thm:complete-four}
For finite $\Sigma$:
$\Sigma\dershort S\varphi \iff \Sigma\models_{\FOUR} S\varphi$.
[L:\,\lean{derives\_iff\_consequence4}]
\end{theorem}
\begin{proof}[Proof sketch]
\emph{Soundness}: induction on the closability derivation. A closed branch
is unsatisfiable since opposite signs are exclusive; each rule preserves
satisfiability downward by the local soundness equivalences read off the
tables (e.g.\ the first bit of $w(\varphi\wedge\psi)$ is $1$ iff both
conjuncts' first bits are). \emph{Completeness}: strong induction on the
total size of undecomposed formulas in the branch. Compound members are
decomposed by their rule (semantic equivalences justify that
unsatisfiability passes to the children); when only literals remain, an
unsatisfiable literal branch must contain an opposite-signed pair --- else
the valuation reading the positive literals off the branch (the canonical
valuation) satisfies it.
\end{proof}

\begin{theorem}[Lifting; continuous completeness]\label{thm:complete-cont}
A finite branch is satisfiable continuously iff satisfiable over $\FOUR$.
Hence for finite $\Sigma$:
$\Sigma\dershort S\varphi \iff \Sigma\models S\varphi$.
[L:\,\lean{satisfiable\_iff\_four}, \lean{derives\_iff\_consequenceC}]
\end{theorem}
\begin{proof}[Proof sketch]
Left to right: project the model's valuation through $\proj{\tau}$;
satisfaction transfers by \cref{thm:projection}. Right to left: embed a
$\FOUR$-valuation at the corners with any threshold (cf.\
\cref{prop:tau}); the projection of the embedded model is the original
valuation, so satisfaction transfers back.
\end{proof}

\begin{theorem}[Decidability]\label{thm:decidable}
For finite $\Sigma$, $\Sigma\dershort S\varphi$ (equivalently
$\Sigma\models S\varphi$) is decidable.
[L:\,\lean{consequence4Bool\_correct}, \lean{decidableDerives},
\lean{decidableConsequenceC}]
\end{theorem}
\begin{proof}
By \cref{thm:complete-four,thm:complete-cont} it suffices to decide
$\models_{\FOUR}$. Evaluation depends only on the atoms occurring in a
formula (structural induction --- the verified lemma), and the occurring
atoms $A$ of $\Sigma\cup\{S\varphi\}$ are finitely many. Hence
$\Sigma\models_{\FOUR}S\varphi$ holds iff it holds for the $4^{|A|}$
valuations extending assignments $A\to\FOUR$ by $\vN$ elsewhere --- a
finite computation on the tables. (Alternatively: fair tableau expansion
terminates, by the usual subformula/finiteness argument, and decides
$\dershort$.)
\end{proof}

\begin{proposition}[Finite complexity]\label{prop:complexity}
\textnormal{[paper]}
Finite signed FOUR consequence is coNP-complete.
\end{proposition}
\begin{proof}
Non-consequence has a polynomial certificate: assign two bits to each atom
occurring in $\Sigma\cup\{S\varphi\}$, evaluate bottom-up, and check that
all premises hold while $S\varphi$ fails. For hardness, reduce Boolean
tautology. Given a Boolean formula $\theta$, read it as an $\hmz$-free NPL
formula $\theta^*$ and add, for every variable $p_i$, the premises
$\Tp(p_i\vee\lnot p_i)$ and $\Tm(p_i\wedge\lnot p_i)$. These force $p_i$
to be $\vT$ or $\vF$; $\vB$ violates the second premise and $\vN$ violates
the first. By classical-corner recovery, $\Gamma_\theta\models_{\FOUR}
\Tp\theta^*$ iff $\theta$ is a Boolean tautology. Boolean tautology is
coNP-complete by the Cook--Karp theory
\cite{cook1971complexity,karp1972reducibility}.
\end{proof}

\begin{theorem}[Compactness; strong completeness]\label{thm:compactness}
\textnormal{[L:\,\lean{compactness\_satisfiable4\_set},
\lean{compactness\_consequence4\_set},
\lean{compactness\_consequenceC\_set},
\lean{derivesSet\_iff\_consequenceCSet}]}
(i) An arbitrary set $\Sigma$ of signed formulas is $\FOUR$-satisfiable iff
every finite subset is. (ii) The same holds for continuous satisfiability.
(iii) Reading $\Sigma\dershort S\varphi$ for infinite $\Sigma$ as ``some
finite subset derives $S\varphi$'':
$\Sigma\dershort S\varphi \iff \Sigma\models S\varphi$ for arbitrary
$\Sigma$.
\end{theorem}
\begin{proof}
(i) One direction is restriction. For the converse, fix the enumeration
$p_0,p_1,\dots$ of atoms (\cref{def:lang}). A \emph{level-$n$ assignment}
is $\sigma:\{p_0,\dots,p_{n-1}\}\to\FOUR$; call $\sigma$ \emph{good} iff
for every finite $\Sigma_0\subseteq\Sigma$ some $\FOUR$-valuation extending
$\sigma$ satisfies $\Sigma_0$. The empty assignment is good (hypothesis).
If $\sigma$ is good, some one-point extension $\sigma\cup\{p_n\mapsto c\}$
is good: otherwise, for each of the four values $c$ pick a finite
$\Sigma_c$ witnessing failure; $\Sigma^*=\bigcup_c\Sigma_c$ is finite, so
some $v\supseteq\sigma$ satisfies $\Sigma^*$; putting $c:=v(p_n)$, $v$
extends $\sigma\cup\{p_n\mapsto c\}$ and satisfies
$\Sigma_c\subseteq\Sigma^*$ --- contradiction. (The case split over $c$ is
exhaustive: $\FOUR$ has exactly four elements.) By recursion along
$\mathbb{N}$ choose a chain $\sigma_0\subseteq\sigma_1\subseteq\cdots$ of
good assignments (countable dependent choice; we work classically). Define
$v^*(p_n):=\sigma_{n+1}(p_n)$. For $S\varphi\in\Sigma$: the atoms of
$\varphi$ lie among $p_0,\dots,p_{n-1}$ for some $n$; goodness of
$\sigma_n$ with $\Sigma_0=\{S\varphi\}$ gives $w\supseteq\sigma_n$
satisfying $S\varphi$; $w$ and $v^*$ agree on the atoms of $\varphi$, so
by finite dependence (\cref{thm:decidable}'s lemma) they give $\varphi$
the same value, and signed satisfaction reads only that value. So $v^*$
satisfies $\Sigma$.
(ii) Both transfer directions of \cref{thm:complete-cont} are memberwise
and never use finiteness of the branch: a continuous model's projected
valuation satisfies each member, and an embedded $\FOUR$-valuation's model
satisfies each member.
(iii) If some finite $\Sigma_0\dershort S\varphi$ then
$\Sigma_0\models S\varphi$ (\cref{thm:complete-cont}), hence
$\Sigma\models S\varphi$ (shrinking premises weakens the hypothesis of the
universally quantified implication). Conversely if no finite subset
derives, then by \cref{thm:complete-four} no finite
$\Sigma_0\models_{\FOUR}S\varphi$, so every finite subset of
$\Sigma\cup\{\bar S\varphi\}$ is satisfiable (failing $S\varphi$ is
satisfying $\bar S\varphi$); by (i)--(ii) the whole of
$\Sigma\cup\{\bar S\varphi\}$ is, so $\Sigma\not\models S\varphi$.
\end{proof}

\section{The behavioural profile of NPL}\label{sec:profile}

\begin{corollary}[Conservativity; nontriviality]\label{cor:conservative}
Unsigned NPL consequence is the all-$\Tp$ fragment of signed consequence
(definitional), and on $\hmz$-free formulas it coincides with Belnap--Dunn
FDE consequence: identical value space, operations, and designated set
define the same matrix consequence \textnormal{[paper]}. Nontriviality:
$\emptyset\not\dershort\Tp p$, by the constant-$\vN$ countermodel
[L:\,\lean{consistency\_witness}]; note every formula evaluates to $\vN$
under the constant-$\vN$ valuation, so NPL has no theorems --- like FDE.
\end{corollary}

\begin{corollary}[Paraconsistency]\label{cor:nonexplosion}
$\{\Tp p,\ \Tp\lnot p\}\not\models \Tp q$: at $p\mapsto\vB$, $q\mapsto\vN$
both premises hold ($t(p)=f(p)=1$) and the conclusion fails ($t(q)=0$).
Explosion, ex falso, and disjunctive syllogism fail.
[L:\,\lean{non\_explosion}]
\end{corollary}

\begin{proposition}[The material conditional does not detach]\label{prop:mp-fails}
$\{\Tp p,\ \Tp(p\Rightarrow q)\}\not\models\Tp q$: same witness,
$t(\lnot p\vee q)=\max(f(p),t(q))=\max(1,0)=1$ yet $t(q)=0$. Consequently
$\vee$-elimination must be taken as a meta-rule (reasoning by cases), not
routed through $\Rightarrow$.
[L:\,\lean{modus\_ponens\_fails}, \lean{disj\_elim}]
\end{proposition}

\begin{proposition}[Derived rules; ND incompleteness]\label{prop:nd}
Let $\ndder$ be the calculus generated by exactly the following rules
(finite premise lists; the source system's rule set):
assumption; $\wedge$-introduction and both eliminations; both
$\vee$-introductions; $\vee$-elimination by cases (from
$\Gamma\vdash\varphi\vee\psi$, $\varphi,\Gamma\vdash\chi$,
$\psi,\Gamma\vdash\chi$ infer $\Gamma\vdash\chi$); $\lnot\lnot$-introduction
and elimination; the four De~Morgan directions for $\wedge$ and $\vee$;
$\hmz$-introduction and both $\hmz$-eliminations. Then:
(a) every $\ndder$-rule is a derived rule of $\dershort$ (each is
semantically valid --- the $\hmz$-rules because the truth channel of $\hmz$
is $\min$ --- and \cref{thm:complete-four} packages validity as
derivability); but
(b) $\ndder$ is \emph{incomplete}: $\hmz$-self-duality
$\lnot(p\hmz q)\models\lnot p\hmz\lnot q$ (a value identity,
\cref{lem:structure}(iii)) is not $\ndder$-derivable.
For (b): no rule of $\ndder$ constrains the falsity channel of $\hmz$ ---
on the truth channel $\hmz$ and $\wedge$ agree --- so reading $\hmz$ as
$\wedge$ throughout preserves the soundness of every rule; but the target's
$\wedge$-reading, $\lnot(p\wedge q)\models\lnot p\wedge\lnot q$, fails at
$p\mapsto\vT$, $q\mapsto\vF$. Self-duality is a falsity-channel fact, and
the four-signed calculus is canonical precisely because its
$F^{\pm}\hmz$ rules reach that channel.
[L:\,\lean{ProofTheory.ND} (16 constructors), \lean{ND.sound\_w} (generic
truth-channel soundness), \lean{nd\_sound}, \lean{nd\_incomplete}]
\end{proposition}

\begin{proposition}[$\hmz$ is not definable from $\{\lnot,\wedge,\vee\}$]\label{prop:primitive}
The $\{\lnot,\wedge,\vee\}$ fragment is classically closed: a $\hmz$-free
formula takes a classical value ($\vT$/$\vF$) whenever its atoms do
(structural induction on the corner tables). Since
$\vT\hmz\vF=\vN$ is not classical, no $\hmz$-free formula defines $\hmz$:
harmonization strictly extends the expressive power of the FDE fragment.
[L:\,\lean{classical\_closed}, \lean{oplus\_not\_definable}]
\end{proposition}

\begin{proposition}[Exact role of $\hmz$; non-congruence]\label{prop:oplus}
$\Tp(\varphi\hmz\psi)$ and $\Tp(\varphi\wedge\psi)$ coincide in every model
(both truth channels are $\min$), so the two formulas are interderivable;
yet at $\varphi\mapsto\vB,\psi\mapsto\vT$: $\lnot(\varphi\wedge\psi)$ is
designated ($\lnot\vB=\vB$) while $\lnot(\varphi\hmz\psi)$ is not
($\lnot\vT=\vF$). Interderivability does not license replacement under
$\lnot$; only value equivalence does. NPL is non-congruential by design.
[L:\,\lean{conj\_oplus\_interderivable}, \lean{oplus\_conj\_Fpos\_fails},
\lean{noncongruence\_witness}]
\end{proposition}

\begin{corollary}[Latent collapse, thresholded]\label{cor:latent}
$V_M(\varphi\hmz\lnot\varphi)=(m,m)$ with $m=\min(t,f)$
(\cref{lem:structure}(ii)); for $\tau>m$ the harmonized contradiction is
unmanifest ($\vN$), not explosive. [L:\,\lean{latent\_collapse}]
\end{corollary}

\section{The generative tier}\label{sec:generative}

Nothing in \cref{sec:syntax,sec:semantics,sec:projection,sec:calculus,%
sec:metatheory,sec:profile} depends on how atoms acquire truth-objects; in
the artifact this is a checkable property of the import graph (the module
\lean{Generative} is imported by none of the core modules). This section
adds the initialization theory.

\begin{definition}[Generative frame; channels; quasivance]\label{def:genframe}
A \emph{generative frame} is $F=(d,\Phi)$ with $d\ge 1$ and
$\Phi:[0,1]^d\to[0,1]$ satisfying $\Phi(1{-}\Theta)=\Phi(\Theta)$ and
$\Phi(\tfrac12,\dots,\tfrac12)=1$. A \emph{support channel}
$c=(\alpha,\Theta)$ has \emph{existence intensity} $\alpha\in[0,1]$ and
\emph{polarization structure} $\Theta\in[0,1]^d$; its effective intensity
is $\alpha\cdot\Phi(\Theta)$. Two independent channels per atom
(for/against) induce the truth-object
$(\alpha_T\Phi(\Theta_T),\,\alpha_F\Phi(\Theta_F))\in[0,1]^2$ (interface
theorem: products of $[0,1]$-values stay in $[0,1]$), so the core applies
to any admissible $\Phi$; the canonical instance
$\Phi_c(\Theta)=(\prod_k(1-2|\Theta_k-\tfrac12|))^{1/d}$ is verified
admissible. A channel is \emph{quasivant} iff $\alpha=0$ with non-neutral
$\Theta$: unmanifest yet structured --- the notion that motivated the
system (structure at probability zero is invisible to entropy-based
measures \cite{shannon1948mathematical}).
[L:\,\lean{GenFrame}, \lean{GenState.init\_mem}, \lean{canonFrame},
\lean{Channel.Quasivant}]
\end{definition}

\begin{proposition}[Quasivance at the interface]\label{prop:quasivance}
(i) A quasivant state initializes to $(0,0)$ and projects to $\vN$ for
every $\tau\in(0,1]$. (ii) \emph{Forgetfulness}: initialization is not
injective; quasivant and non-quasivant states can initialize identically,
so quasivance is not recoverable from the truth-object --- claims about
$\vN$ in the core concern unmanifestness only. (iii) In the canonical
frame, one fully polarized component annihilates effective intensity even
at $\alpha=1$. (iv) Initialization is flip-covariant and surjective onto
the square. We state (ii) prominently: it is the honest price of the
generative reading.
[L:\,\lean{quasivant\_projects\_N}, \lean{init\_not\_injective},
\lean{polar\_kills\_intensity}, \lean{init\_flip}, \lean{init\_surjective}]
\end{proposition}

\section{Related work}\label{sec:related}

\textbf{Belnap--Dunn / FDE.} The $\FOUR$ fragment on $\{\lnot,\wedge,\vee\}$
\emph{is} FDE \cite{belnap1977useful,belnap1977computer,dunn1976intuitive};
see \cite{omori2017forty} for the landscape of FDE extensions, within which
NPL sits as a consensus expansion with a projection layer. \textbf{Logical
bilattices.} The square is the product bilattice
\cite{ginsberg1988multivalued,fitting1991bilattices}; $\hmz$ is consensus;
complete Gentzen/Hilbert calculi for the four-valued bilattice language are
due to Arieli--Avron \cite{arieli1996reasoning} (see also
\cite{rivieccio2010algebraic}). The precise relation:

\begin{proposition}[Position within the bilattice framework]\label{prop:bifilter}
\textnormal{[paper]}
(i) For every $\tau\in(0,1]$, $D_\tau=\{(t,f):t\ge\tau\}$ is a prime
bifilter of $[0,1]\odot[0,1]$ (the five defining conditions hold
componentwise on the truth channel), so
$\langle[0,1]\odot[0,1],D_\tau\rangle$ is a logical bilattice; by the
Arieli--Avron collapse theorem (every logical bilattice determines the same
consequence as $\langle\FOUR,\{t,\top\}\rangle$ ---
\cite[Thm.~2.17]{arieli1996reasoning}, as presented in
\cite[Thm.~2.1.4]{rivieccio2010algebraic}), the \emph{unsigned} NPL
consequence is exactly the $\{\wedge,\vee,\otimes,\lnot\}$-fragment of
their logic $\mathrm{LB}$ --- subsumed, and consistent with our independent
\cref{thm:complete-cont}. (ii) The satisfaction set of $\Tm$,
$\{(t,f):t<\tau\}$, contains $(0,0)$ but not $(1,1)$, hence is not upward
closed in $\lek$ and is not a bifilter (bifilters are upward closed in both
orders and contain the knowledge top). The four-signed relation is
therefore not a logical-bilattice matrix consequence; the signed layer is
where NPL leaves the framework. (iii) Nor can the negative signs be
internalized: no formula $\psi$ satisfies
$\Tp\psi\iff\Tm p$ in all valuations, since the constant-$\vB$ valuation is
a fixpoint of every connective --- every formula is designated there, while
$\Tm p$ fails. [L:\,\lean{signs\_not\_internalizable}, via
\lean{eval\_const\_B}]
\end{proposition}

\textbf{Logic programming.} Fitting \cite{fitting1991bilattices}
interprets logic programs over continuous bilattices with the knowledge
operations (including consensus) available in rule bodies; that line
provides fixed-point semantics for programs, not a consequence relation
with a completeness theorem for a propositional language --- the two uses
of the same algebra are complementary. \textbf{Paraconsistent G\"odel
logic.} The continuous $\{\lnot,\wedge,\vee\}$ clauses coincide with the
propositional base of \cite{bilkova2022paraconsistent} (twist product of
$[0,1]$); their language has a G\"odel implication and no consensus, ours
has consensus and no implication, and the consequence relations differ
(signed vs.\ order-based). \textbf{Fuzzy sets.} The projection core is
$\alpha$-cut cutworthiness \cite{klir1995fuzzy}. \textbf{Tableaux.}
Sets-as-signs tableaux are standard \cite{haehnle1994automated}; FDE has
textbook signed tableaux \cite{priest2008introduction}. The closest tableau
precedent for our setting is the \emph{constraint tableaux} of
\cite{bilkova2021constraint} for two-dimensional \L{}ukasiewicz and
G\"odel logics over the same twist product of $[0,1]$: there, branches
carry order constraints on real values because the fuzzy implications do
not reduce to any finite matrix. Our calculus needs only four \emph{finite}
signs precisely because NPL has no implication --- all connectives are
cut-worthy, exact projection holds, and consequence collapses onto $\FOUR$
(\cref{thm:projection}); adding a G\"odel implication to NPL would break
that theorem. The two designs are complementary solutions for different
languages. \textbf{LP.} Designating $\vB$ descends from
\cite{priest1979logic}. \textbf{Verified metatheory.} Paraconsistent logics
have been formalized in proof assistants before --- notably an
infinite-valued paraconsistent logic with its metatheory in Isabelle
\cite{villadsen2017formalizing}; we are not aware of verified
\emph{completeness} for a consensus-bearing logic over $[0,1]^2$.

To our knowledge --- search scope documented in the project's related-work
audit; last sweep July 2026, covering post-2022 twist-product,
bilattice-calculus, and proof-assistant literature --- no prior system
combines a consensus connective over $[0,1]^2$ with four-signed consequence
and a machine-checked completeness theorem.

\section{Conclusion and open problems}\label{sec:conclusion}

NPL's propositional development is closed end to end: generative
initialization, continuous two-channel semantics, exact projection, a
four-signed calculus, soundness/completeness/decidability verified in
Lean, compactness, strong completeness, classical-corner recovery verified in Lean, coNP-complete finite consequence, a provably
incomplete ND fragment explaining why the signed calculus is canonical,
and a precise statement of which parts of the system are and are not
instances of the logical-bilattice framework. Since this manuscript snapshot,
the project has also Lean-verified completeness for the $\hmz$-De-Morgan
extension of $\ndder$, and has closed the finite-domain quantified program:
a function-free quantified extension over finite nonempty domains with crisp
equality, an exact threshold projection, and a machine-checked semantic
completeness theorem for a macro-free extensional tableau calculus (a
companion manuscript is in preparation in the repository under
\texttt{papers/npl-finite-fo/}). Open:
the full infinite-domain quantified extension; the emergence layer of the source material
remains outside the formal development by design.

\appendix
\section{The Lean artifact}\label{app:lean}
\textbf{Availability.} The \texttt{Nullivance} library will be archived
with a DOI at submission; until then it is available from the author on
request. \textbf{Toolchain.} Lean 4 (\texttt{leanprover/lean4:v4.32.0-rc1})
with mathlib pinned via the Lake manifest; build with
\texttt{lake exe cache get \&\& lake build} (approx.\ 2000 jobs; the
library itself compiles in under a minute on commodity hardware).
\textbf{Coverage policy.} The build is \texttt{sorry}-free. Every result
marked [L] is verified as stated, with the following systematic
correspondences: branches are lists (finite sets up to the verified
monotonicity lemma); ``some tableau closes'' is the inductive predicate
\lean{Closes} (two closure, sixteen rule constructors); the calculus-level
results are stated against \lean{Closes} directly; explicit finite tableau
proof trees are equivalent to \lean{Closes}. Results marked [paper]
--- \cref{cor:conservative} (FDE-conservativity half), \cref{prop:complexity},
\cref{prop:bifilter} ---
carry full proofs in this paper and are not formalized; the termination and
fair-search notions appear only in the alternative proof remark of
\cref{thm:decidable}. \textbf{Modules.} \texttt{Syntax}, \texttt{Semantics}
($\FOUR$), \texttt{Continuous} (square, projection, bilattice orders),
\texttt{ProofTheory} (\lean{Closes}, \lean{ND}, \lean{NDO}), \texttt{Tableau}
(explicit proof trees), \texttt{Metatheory}
(soundness, completeness, lifting, $\tau$-invariance, behavioural profile),
\texttt{Decidability} (finite model checker and \lean{Decidable} instances),
\texttt{Compactness} (Set/Finset API, compactness, strong completeness),
\texttt{Classical} (Boolean recovery on the glut/gap-free fragment),
\texttt{FiniteFO} (the complete finite-domain quantified layer, including its
verified completeness theorem — companion manuscript),
\texttt{Generative} (\cref{sec:generative}; imported by no core module).

\bibliographystyle{plain}
\bibliography{../../references/bibliography}

\end{document}
